\section{Preliminaries and Problem Formulation}
\label{sec:problem}

\subsection{Hierarchical Edge Learning System}
\label{subsec:system}

We consider a cloud-edge-client learning system with one cloud server, a set of edge servers
$\mathcal{E}=\{1,\ldots,E\}$, and a set of clients $\mathcal{C}=\{1,\ldots,C\}$. Each client
$i$ is associated with one edge server $e(i)\in\mathcal{E}$. Client $i$ owns a local dataset
$\mathcal{D}_i=\mathcal{D}_i^l\cup\mathcal{D}_i^u$, where $\mathcal{D}_i^l$ and
$\mathcal{D}_i^u$ denote labeled and unlabeled samples, respectively. The labeled ratio is
usually small in practical edge intelligence tasks because data annotation is expensive, while
unlabeled data are continuously generated by distributed devices.

Unlike conventional federated learning where each client directly trains a local model, we
target resource-constrained edge networks in which clients are lightweight data participants.
The major training workload is moved from clients to edge servers. Each selected client uploads
its mini-batch data or features to its associated edge server, and the edge server performs local
model updates. The cloud server coordinates cross-edge knowledge sharing. This design matches
edge environments where clients have limited computation, energy, and wireless connectivity,
while edge servers and the cloud have stronger computation and more stable communication links.

Let $f_\phi(\cdot)$ denote a shared feature encoder and $h_{\theta_e}(\cdot)$ denote the
classifier maintained by edge server $e$. In this work, the encoder is used as a feature
extractor and the edge classifier is updated during training. The cloud maintains a teacher
classifier $h_{\theta_c}(\cdot)$, which is updated from edge models through cloud-side
aggregation.

\subsection{Learning Objective Under Network Constraints}
\label{subsec:objective}

The goal is to train an accurate model under data heterogeneity, label scarcity, bandwidth
constraints, and client instability. At each communication round $t$, the system selects a
subset of clients $\mathcal{S}_t\subseteq\mathcal{C}$ and trains edge models with their data.
The client selection must satisfy a round communication budget:
\begin{equation}
    \sum_{i\in\mathcal{S}_t} m_i^{\mathrm{ce}}
    + \sum_{e\in\mathcal{E}_t} m_e^{\mathrm{ec}}
    + \sum_{e\in\mathcal{E}_t} m_e^{\mathrm{ceg}}
    \le B_t,
    \label{eq:budget}
\end{equation}
where $m_i^{\mathrm{ce}}$ is the client-to-edge upload size, $m_e^{\mathrm{ec}}$ is the
edge-to-cloud upload size, $m_e^{\mathrm{ceg}}$ is the cloud-to-edge downlink size, and $B_t$
is the communication budget of round $t$. $\mathcal{E}_t$ denotes the set of active edge
servers in round $t$.

We also model the round latency as a three-level cloud-edge-client process:
\begin{equation}
    T_t = T_t^{\mathrm{c2e}}
        + T_t^{\mathrm{edge}}
        + T_t^{\mathrm{e2c}}
        + T_t^{\mathrm{cloud}}
        + T_t^{\mathrm{c2e\_down}},
    \label{eq:latency}
\end{equation}
where $T_t^{\mathrm{c2e}}$ is the terminal-to-edge upload latency, $T_t^{\mathrm{edge}}$ is
the edge computation time, $T_t^{\mathrm{e2c}}$ is the edge-to-cloud upload latency,
$T_t^{\mathrm{cloud}}$ is the cloud aggregation time, and $T_t^{\mathrm{c2e\_down}}$ is the
cloud-to-edge downlink latency. This formulation allows us to evaluate not only accuracy but
also communication cost and time-to-target-accuracy.

\section{The Proposed EdgeQGFed Framework}
\label{sec:method}

EdgeQGFed is a communication-efficient cloud-edge-client collaborative learning framework. It
contains three key components: edge-side semi-supervised training, cloud-side quality-aware graph
attention aggregation, and resource-aware client sampling under network constraints.

\subsection{Edge-Side Semi-Supervised Training}
\label{subsec:edge_training}

At round $t$, each active edge server $e$ receives mini-batches from selected active clients
under its coverage. For an input $x$, the encoder produces a representation
$z=f_\phi(x)$, and the edge classifier outputs logits $h_{\theta_e}(z)$. For labeled samples,
the edge server minimizes the supervised cross-entropy loss:
\begin{equation}
    \mathcal{L}_{\mathrm{sup}}^e =
    \frac{1}{|\mathcal{B}_e^l|}
    \sum_{(x,y)\in\mathcal{B}_e^l}
    \ell_{\mathrm{ce}}\left(h_{\theta_e}(f_\phi(x)), y\right),
    \label{eq:sup_loss}
\end{equation}
where $\mathcal{B}_e^l$ is the labeled batch at edge $e$.

For unlabeled samples, EdgeQGFed uses cloud-edge consistency. Given a weakly augmented input
$x^w$ and a strongly augmented input $x^s$, the edge prediction and cloud prediction are:
\begin{equation}
    p_e = \mathrm{softmax}\left(h_{\theta_e}(f_\phi(x^s))/T_p\right), \quad
    p_c = \mathrm{softmax}\left(h_{\theta_c}(f_\phi(x^w))/T_p\right),
    \label{eq:edge_cloud_pred}
\end{equation}
where $T_p$ is the pseudo-label temperature. The consistency loss is:
\begin{equation}
    \mathcal{L}_{\mathrm{con}}^e =
    d(p_e, p_c),
    \label{eq:consistency_loss}
\end{equation}
where $d(\cdot,\cdot)$ can be mean squared error, $\ell_1$ distance, or KL divergence.

To reduce confirmation bias, pseudo-labels are weighted by both edge confidence and cloud
confidence. Let $\hat{y}_e=\arg\max p_e$, $\hat{y}_c=\arg\max p_c$, $q_e=\max p_e$, and
$q_c=\max p_c$. A sample is used for pseudo-label training only when the edge and cloud
predictions agree. Its soft weight is:
\begin{equation}
    w(x) =
    \mathbf{1}[\hat{y}_e=\hat{y}_c]
    \cdot
    \sigma\left(\frac{q_e-\tau_e}{\gamma}\right)
    \cdot
    \sigma\left(\frac{q_c-\tau_c}{\gamma}\right),
    \label{eq:pseudo_weight}
\end{equation}
where $\tau_e$ and $\tau_c$ are edge and cloud confidence thresholds, $\gamma$ is a temperature
for weight smoothing, and $\sigma(\cdot)$ is the sigmoid function. The weighted pseudo-label
loss is:
\begin{equation}
    \mathcal{L}_{\mathrm{pl}}^e =
    \frac{\sum_{x\in\mathcal{B}_e^u} w(x)
    \ell_{\mathrm{ce}}\left(h_{\theta_e}(f_\phi(x^s)), \hat{y}_c\right)}
    {\sum_{x\in\mathcal{B}_e^u} w(x)+\epsilon}.
    \label{eq:pseudo_loss}
\end{equation}

The overall edge training objective is:
\begin{equation}
    \mathcal{L}_e =
    \mathcal{L}_{\mathrm{sup}}^e
    + \lambda_t \mathcal{L}_{\mathrm{con}}^e
    + \lambda_p \mathcal{L}_{\mathrm{pl}}^e,
    \label{eq:edge_loss}
\end{equation}
where $\lambda_t$ is gradually increased by a ramp-up schedule and $\lambda_p$ controls the
pseudo-label loss.

\subsection{Edge Summary Construction}
\label{subsec:summary}

After local edge training, each active edge server uploads a compact summary to the cloud. The
summary contains model parameters, class prototypes, pseudo-label quality, labeled ratio,
prediction confidence, and batch size:
\begin{equation}
    s_e = \{\theta_e, P_e, n_e, \rho_e, r_e, c_e, b_e\}.
    \label{eq:summary}
\end{equation}
Here $P_e$ is the class prototype matrix, $n_e$ is the class-wise prototype count, $\rho_e$ is
the average pseudo-label weight, $r_e$ is the labeled ratio, $c_e$ is the mean cloud confidence,
and $b_e$ is the number of processed samples. For class $k$, the prototype is computed as:
\begin{equation}
    P_e^k =
    \frac{1}{|\mathcal{B}_{e,k}^l|}
    \sum_{(x,y)\in\mathcal{B}_{e,k}^l} f_\phi(x),
    \label{eq:prototype}
\end{equation}
where $\mathcal{B}_{e,k}^l$ is the labeled subset of class $k$ at edge $e$.

\subsection{Cloud-Side Quality-Aware Graph Attention Aggregation}
\label{subsec:graph_aggregation}

The cloud is not treated as a passive parameter averager. Instead, it constructs a
quality-aware graph over active edge servers and performs graph attention aggregation. For edge
$e$, we first compute a model signature $u_e$ from its parameters by concatenating layer-wise
statistics such as mean, standard deviation, and normalized norm. We also compute a prototype
signature $v_e$ by flattening and normalizing the count-weighted prototype matrix.

For two active edges $i$ and $j$, the cloud computes their attention score as:
\begin{align}
    a_{ij} =
    &\ \mathrm{cos}(u_i,u_j)
    + \alpha_p \mathrm{cos}(v_i,v_j)
    + \alpha_q (1-|\rho_i-\rho_j|) \nonumber \\
    &+ \alpha_r (1-|r_i-r_j|)
    + \alpha_c (1-|c_i-c_j|)
    + \beta \mathbf{1}[i=j],
    \label{eq:attention_score}
\end{align}
where $\alpha_p$, $\alpha_q$, $\alpha_r$, and $\alpha_c$ control the contributions of prototype
similarity, pseudo-label reliability, labeled ratio similarity, and confidence similarity.
$\beta$ is a self-bias term that preserves edge-specific knowledge.

The attention matrix is obtained by temperature-scaled softmax:
\begin{equation}
    A_{ij} =
    \frac{\exp(a_{ij}/\tau_g)}
    {\sum_{k\in\mathcal{E}_t}\exp(a_{ik}/\tau_g)},
    \label{eq:attention_matrix}
\end{equation}
where $\tau_g$ is the graph temperature. A small lower bound is applied to avoid zero attention,
followed by row-wise renormalization.

The cloud generates personalized parameters for each edge:
\begin{equation}
    \tilde{\theta}_i =
    \sum_{j\in\mathcal{E}_t} A_{ij}\theta_j.
    \label{eq:personalized_param}
\end{equation}
Compared with FedAvg, this aggregation allows edge $i$ to receive more knowledge from reliable
and semantically similar edges, while reducing the influence of low-quality or mismatched
updates.

The global cloud parameter is then computed by reliability-aware weighted aggregation:
\begin{equation}
    \bar{\theta} =
    \sum_{i\in\mathcal{E}_t}
    \omega_i \tilde{\theta}_i,
    \quad
    \omega_i =
    \frac{b_i(\rho_i+1)}
    {\sum_{j\in\mathcal{E}_t} b_j(\rho_j+1)}.
    \label{eq:global_param}
\end{equation}

Finally, the cloud teacher is updated by adaptive exponential moving average:
\begin{equation}
    \theta_c \leftarrow
    \mu \theta_c + (1-\mu)\bar{\theta},
    \label{eq:ema}
\end{equation}
where $\mu$ is the EMA decay. The personalized parameters $\tilde{\theta}_i$ are sent back to
edge servers for the next round.

\subsection{Resource-Aware Client Sampling}
\label{subsec:sampling}

To improve communication efficiency, EdgeQGFed selects clients according to both data utility
and network conditions. For client $i$, we define a resource-aware utility score:
\begin{equation}
    U_i =
    \frac{
    \eta_l l_i + \eta_d d_i + \eta_b b_i^{\mathrm{up}} + \eta_a a_i
    }
    {(T_i^{\mathrm{up}})^{\eta_t}},
    \label{eq:client_utility}
\end{equation}
where $l_i$ is the labeled ratio, $d_i$ is the normalized data size,
$b_i^{\mathrm{up}}$ is the normalized uplink bandwidth, $a_i=1-p_i^{\mathrm{drop}}$ is the
availability score, and $T_i^{\mathrm{up}}$ is the estimated upload latency. The coefficients
$\eta_l$, $\eta_d$, $\eta_b$, $\eta_a$, and $\eta_t$ control the importance of data quality,
data quantity, bandwidth, availability, and latency cost.

At each round, EdgeQGFed selects high-utility clients while respecting the maximum number of
parallel uploads per edge and the round communication budget in Eq.~\eqref{eq:budget}. This
sampling strategy avoids wasting communication resources on clients with low data value, poor
links, or high dropout probability.

\subsection{Overall Training Procedure}
\label{subsec:procedure}

The overall training process is as follows. In each round, the system first performs
resource-aware client sampling under communication constraints. Active clients upload data or
features to their associated edge servers. Each edge server updates its classifier using
supervised loss, cloud-edge consistency loss, and confidence-aware pseudo-label loss. Then,
active edges upload summaries to the cloud. The cloud constructs a quality-aware graph, computes
attention-based personalized parameters, updates the cloud teacher by EMA, and sends
personalized parameters back to edges. This loop continues until the total number of training
rounds is reached.

By combining hierarchical workload placement, graph-based cloud coordination, and
network-aware sampling, EdgeQGFed targets the key challenge of edge intelligence: reaching high
accuracy under limited bandwidth, heterogeneous links, label scarcity, and unstable client
participation.
